\label{anlage:typologie}

Die Tabellen dienen als Arbeitsgrundlage für tragende und tragwerksnahe Bauteile, die erneut eingesetzt werden sollen. Entscheidend ist nicht die Bezeichnung aus dem Herkunftsbauwerk, sondern die Form und der konstruktive Zusammenhang, die nach Ausbau oder Zuschnitt tatsächlich erhalten bleiben.

\begin{Table}[title=Grunddaten zum Bauteil, key=tab:typologie-grunddaten]
\begin{SemioTable}{@{}T{0.22\linewidth-2\tabcolsep}T{0.78\linewidth-2\tabcolsep}@{}}
\SemioTableHeaderRow{Bereich & Erfassung}
\SemioTableRow{Herkunft & Bauteil-ID, Quellbauwerk, ursprünglicher Einbauort, Baualter, Stückzahl und Verfügbarkeit}
\SemioTableRow{Geometrie & Aufmaß, Orientierung, Masse, Verformungen, Toleranzen sowie vorhandene Schnitt- und Handhabungspunkte}
\SemioTableRow{Material und Aufbau & Materialart, relevante Kennwerte und konstruktive Details; bei Stahlbeton zum Beispiel Bewehrung, Vorspannung, Betondeckung und Verankerung}
\SemioTableRow{Zustand und Vorgeschichte & Schäden, Exposition, frühere Lagerung und Anschlüsse, bekannte Lastgeschichte sowie Ausbau-, Schnitt- und Reparaturspuren}
\SemioTableRow{Nachweisstand & Herkunft der Angaben, Prüfverfahren, Genauigkeit, Unsicherheiten und offene Punkte}
\end{SemioTable}
\end{Table}

Die folgenden Tabellen ordnen jede Typologie einer der vier Typologiefamilien zu. Für jede Typologie nennen sie Beispiele, die ursprüngliche Tragrolle, naheliegende Wiederverwendungsrollen, fachlich zu prüfende Rollenwechsel und typologiespezifische Zusatzparameter (\tableref{tab:typologien-lineare}, \tableref{tab:typologien-flaechige}, \tableref{tab:typologien-knoten}, \tableref{tab:typologien-rahmen}).

\subsection{Lineare Elemente}

\SemioTableLong[text-size=7.2pt, key=tab:typologien-lineare]{Lineare Elemente}{0.14,0.20,0.16,0.20,0.14,0.16}{Typologie & Beispiele & Ursprüngliche Tragrolle & Naheliegende Wiederverwendungsrollen & Fachlich zu prüfende Rollenwechsel & Typologiespezifische Zusatzparameter}{%
  \SemioTableRow{Träger-\SemioSlash{}Riegelelement & \SemioCellList[,]{Träger, Unterzug, Balken, Riegel, Vollwandbinder, Pfette, Sturz} & Vorwiegend Biegung und Querkraft; je nach System zusätzlich Normalkraft oder Torsion & Biegebeanspruchtes lineares Tragglied, etwa Träger, Riegel, Pfette, Sturz oder Vollwandbinder & \SemioCellList{Stütze oder Pfosten; axiales Stab- oder Aus\-steifungs\-element} & \SemioCellList{Länge und Achsverlauf; Querschnitt und Hauptachsen; Vorkrümmung; Aussparungen und Ausklinkungen; End-, Schnitt- und Anschluss\-bereiche}}
  \SemioTableRow{Stützen-\SemioSlash{}Pfostenelement & \SemioCellList[,]{Stütze, Pfeiler, tragender Pfosten} & Vorwiegend Drucknormalkraft, häufig zusammen mit Biegung & Druckbeanspruchtes vertikales Tragglied, etwa Stütze, Pfeiler oder Pfosten & \SemioCellList{Träger oder Riegel; axiales Stab- oder Aus\-steifungs\-element} & \SemioCellList{Länge und Achsverlauf; Querschnitt; Geradheit und Exzentrizitäten; Kopf-, Fuß-, Schnitt- und Anschluss\-bereiche}}
  \SemioTableRow{Axiales Stab-\SemioSlash{}Aus\-steifungs\-element & \SemioCellList[,]{Zugstab, Druckstab, Fachwerkstab, Diagonale, Verbandsstab} & Zug, Druck oder wechselnde Normalkraft innerhalb eines Stab- oder Aus\-steifungs\-systems & Axiales Stab-, Fachwerk- oder Aussteifungsglied & \SemioCellList{Träger oder Riegel; Stütze oder Pfosten} & \SemioCellList{Stablänge und Querschnitt; Geradheit; Endausbildung; Anschlussversätze; Verankerungen und Kraft\-einleitungs\-bereiche}}
  \SemioTableRow{Flexibles Zugelement -- Sonderform & \SemioCellList[,]{Seil, Litze, Drahtseil, flexibles Zugband} & Zug\-beanspruchung & Abspann-, Hänge- oder zugbeanspruchtes Aus\-steifungs\-element & \SemioCellList{—} & \SemioCellList{Länge und Zugquerschnitt; Aufbau und Dehnsteifigkeit; Vorspannung und Durchhang; Ermüdungs- und Schadensgeschichte; Verankerung und Wieder\-veranker\-barkeit}}
}

\subsection{Flächige Elemente}

\SemioTableLong[text-size=7.2pt, key=tab:typologien-flaechige]{Flächige Elemente}{0.14,0.20,0.16,0.20,0.14,0.16}{Typologie & Beispiele & Ursprüngliche Tragrolle & Naheliegende Wiederverwendungsrollen & Fachlich zu prüfende Rollenwechsel & Typologiespezifische Zusatzparameter}{%
  \SemioTableRow{Plattenelement & \SemioCellList[,]{Deckenplatte, Dachplatte, Podestplatte, Deckenplattenfragment, geschnittener Plattenstreifen} & Lastabtrag überwiegend senkrecht zur Fläche; teilweise zusätzliche Scheibenwirkung & Decken-, Dach-, Podest- oder sonstiges Plattenelement beziehungsweise Plattensegment & \SemioCellList{Wand oder aussteifende Scheibe; nichttragender Raum\-abschluss} & \SemioCellList{Länge, Breite und Dicke; Ebenheit oder Krümmung; Öffnungen; Rest- und Kantengeometrie; Auflager-, Schnitt- und Anschluss\-bereiche}}
  \SemioTableRow{Wand-\SemioSlash{}Scheiben\-element & \SemioCellList[,]{Wandplatte, Wandscheibe, Wandtafel, Wandfragment oder Wandsegment mit oder ohne Öffnung} & Vertikaler Lastabtrag, Scheibenwirkung und\SemioSlash{}oder Raum\-abschluss & Tragwand, aussteifende Scheibe oder nichttragender Raum\-abschluss & \SemioCellList{Decken-, Dach-, Podest- oder sonstige Platten\-funktion} & \SemioCellList{Länge, Höhe und Dicke; Öffnungen; Steg-, Sturz-, Brüstungs- und Randbereiche; Fuß-, Auflager-, Schnitt- und Anschluss\-bereiche}}
}

\subsection{Knoten- und Übergangselemente}

\SemioTableLong[text-size=7.2pt, key=tab:typologien-knoten]{Knoten- und Übergangselemente}{0.14,0.20,0.16,0.20,0.14,0.16}{Typologie & Beispiele & Ursprüngliche Tragrolle & Naheliegende Wiederverwendungsrollen & Fachlich zu prüfende Rollenwechsel & Typologiespezifische Zusatzparameter}{%
  \SemioTableRow{Lokales Knoten-\SemioSlash{}Übergangs\-element & \SemioCellList[,]{Decke--Träger-Fragment, Decke--Stütze-Fragment, Träger--Stütze-Fragment, Decke--Wand-Fragment, Konsolenfragment, Stützenkopf- oder Stützenfußfragment, Auflager- oder Verankerungsmodul} & Lokale Einleitung, Umlenkung oder Weiterleitung von Kräften & Anschluss-, Auflager-, Übergangs-, Verankerungs- oder Lastverteilungselement in einer vergleichbaren Konfiguration & \SemioCellList{—} & \SemioCellList{Restgeometrie; Lage der Teilbereiche zueinander; Kontakt- und Auflagerflächen; konstruktive Kontinuität; Verbindungsmittel; Schnitt- und Anschluss\-bereiche}}
  \SemioTableRow{Mehrgliedriges Knoten-\SemioSlash{}Übergangs\-element & \SemioCellList[,]{Decke--Träger--Stütze-Fragment und ähnliche lokale Konfigurationen} & Lokale Kraftübertragung zwischen mehreren linearen und\SemioSlash{}oder flächigen Teilbereichen & Mehrgliedriger Anschluss oder Rahmenknoten; lokales Auflager- oder Lastverteilungsmodul & \SemioCellList{—} & \SemioCellList{Lokale Topologie; Anzahl und Lage der Teilbereiche; Restlängen und Restabmessungen; Kraftübertragungsflächen; Schnitt- und Anschluss\-bereiche}}
}

\subsection{Rahmen- und Systemelemente}

\SemioTableLong[text-size=7.2pt, key=tab:typologien-rahmen]{Rahmen- und Systemelemente}{0.14,0.20,0.16,0.20,0.14,0.16}{Typologie & Beispiele & Ursprüngliche Tragrolle & Naheliegende Wiederverwendungsrollen & Fachlich zu prüfende Rollenwechsel & Typologiespezifische Zusatzparameter}{%
  \SemioTableRow{Rahmenelement & \SemioCellList[,]{Rahmenfragment, Portalrahmen, Portalrahmensegment, Fertigteilrahmen} & Gemeinsame Tragwirkung mehrerer linearer Glieder und ihrer Knoten & Rahmen-, Portal- oder Aussteifungsrahmensegment beziehungsweise -modul & \SemioCellList{—} & \SemioCellList{Gesamtgeometrie; Spannweiten und Höhen; Teilquerschnitte; Vollständigkeit; Knoten und Verbindungen; Stabilisierung aus der Ebene; Transport- und Montage\-grenzen}}
  \SemioTableRow{Fachwerk-\SemioSlash{}Verbands\-system & \SemioCellList[,]{Fachwerkfragment, Fachwerkfeld, Verbandsfragment} & Systemwirkung aus mehreren überwiegend axial beanspruchten Stäben & Fachwerk-, Verbands- oder Aussteifungssegment beziehungsweise -modul & \SemioCellList{—} & \SemioCellList{Topologie; Stabanordnung und Stablängen; Teilquerschnitte; Systemhöhe; Knoten und Anschlüsse; Vollständigkeit und Stabilisierung aus der Ebene}}
  \SemioTableRow{Trägerrost-\SemioSlash{}Raumrahmen\-system & \SemioCellList[,]{Träger\-rost\-fragment, Raumrahmen\-fragment} & Flächig oder räumlich gekoppelte Tragwirkung mehrerer Glieder und Knoten & Trägerrost-, Raumrahmen- oder räumliches Stabtragwerkssegment beziehungsweise -modul & \SemioCellList{—} & \SemioCellList{Raster und Systemgeometrie; Gliederhierarchie; Teilquerschnitte; Knoten und Verbindungen; Lager- und Stabilisierungspunkte; Transport- und Montage\-grenzen}}
  \SemioTableRow{Gemischte Tragwerks\-baugruppe & \SemioCellList[,]{Größere Decken--Träger--Stützen-Baugruppe oder gemischtes Tragwerksmodul mit erhaltener Systemwirkung} & Gemeinsamer Lastabtrag unterschiedlicher linearer und flächiger Teilbereiche & Tragwerksmodul entsprechend der erhaltenen Systemwirkung & \SemioCellList{—} & \SemioCellList{Gesamtgeometrie und Aufbau; Lage der Komponenten; Vollständigkeit; Systemgrenzen; Knoten und Verbindungen; fehlende Teile; Schnitt-, Transport- und Montage\-grenzen}}
}
